\section{From neuron values instead of outputs (E6--E8)}
\label{sec:wsp}
In E5, Claude built edits behaviorally. It wrote examples, the harness turned them into a dense vector, and Claude tuned a strength by reading outputs. It never looked at a neuron. That is not how Weight-Space Programming worked~\cite{takehi2026wsp}. There, the frontier model had a persistent kernel over the raw weights plus a small budget of forward passes that returned MLP activations and hidden states. It localized the problem to specific MLP output columns from those values and edited only those columns.

We gave Claude the same kind of instruments over Qwen3.5-4B, running on Kaggle (2$\times$T4, the original Hugging Face weights in FP32). Every comparison in this section is regenerated there.
\begin{itemize}
\item \textbf{Python kernel.} A persistent kernel over the weights and over saved reads, with no access to model calls.
\item \textbf{Activation reads.} \texttt{read\_activations} is a teacher-forced pass that returns every MLP neuron's activation, the residual stream and, optionally, a \emph{direct logit attribution} of every neuron to chosen target tokens at the positions where they are emitted.
\item \textbf{Paired generations.} \texttt{generate} runs the original and the edited model side by side.
\item \textbf{Column edits.} \texttt{save\_edit} changes specific \texttt{W\_down} columns, applied as $\Delta W\,x_S$ from the layer's own neuron activations.
\item \textbf{Budget.} Each task gets 10 forward passes.
\end{itemize}
On a tiny model, a column edit is bit-for-bit an explicit weight edit, and the attribution read predicts the real logit drop from silencing a neuron within 5\%.

\subsection{E6: plain-text math from neuron values}
\label{sec:e6}
The task is the one E5's weight edit failed: Qwen should write arithmetic in plain text, not LaTeX, and keep the numbers. Claude had 10 forward passes on practice prompts and used 8. Table~\ref{tab:e6} lists them.

\begin{table}[t]\centering\small
\begin{tabular}{rp{4.9cm}p{5.4cm}p{1.9cm}}\toprule
Pass & What Claude ran & What it found or did & LaTeX per 100 words\\\midrule
1 & 24 normal answers & LaTeX in 24/24 & 20.8\\
2 & read: LaTeX answers vs.\ the same answers with the LaTeX removed & mean-activation contrast plus weight lenses & --\\
3 & silence 26 rule-chosen neurons & no effect & 21.5 (normal 21.2)\\
4 & read with logit attribution to LaTeX openers & L27 n1173 adds +2.2 logits; the next largest adds 0.73 & --\\
5 & silence L27 n1173 alone & no effect & 20.9\\
6 & silence the top 256 by attribution & \texttt{\$\$} blocks become inline \texttt{\$} & 18.2\\
7 & read with attribution that also covers \texttt{\$} and \texttt{\ \$} & the top 512 neurons hold only 26\% of the positive attribution & --\\
8 & silence the top 512 (1.31M weights) & LaTeX in 15/16, numbers intact & 10.9\\\bottomrule
\end{tabular}
\caption{E6's eight practice passes (16 practice prompts each, except pass~1).}
\label{tab:e6}
\end{table}

\paragraph{Localization was precise. The behavior was not local.}
Logit attribution named the neurons that push LaTeX tokens, and it predicts the logit drop from silencing a neuron within 5\% (Section~\ref{sec:wsp}). But silencing even the single largest one, L27 n1173, changed nothing. The model routed around each cut:
\begin{itemize}
\item Removing the writers of \texttt{\$\$} moved display math to inline \texttt{\$}.
\item Removing those as well halved LaTeX density, but almost every answer still had some.
\end{itemize}
This matches what Weight-Space Programming reported for small budgets~\cite{takehi2026wsp}: the attribution is right, but one direction is written by many neurons, and the rest compensate.

\paragraph{Held-out result: H6 fails on its pre-registered metric.}
On E5's 20 held-out arithmetic prompts (seeds 2000--2019, 500 tokens), the 512-neuron edit \texttt{tex-dla512} left LaTeX in 20/20 answers, the same as normal Qwen. The instruction left it in 14/20. All 20 edited answers kept the correct number, as did all 20 normal ones.

Measured as density, which we did not pre-register, the edit cut LaTeX markup from 25.6 to 10.1 tokens per 100 words. The paired difference is $-15.5$ (95\% bootstrap CI $[-18.2,-13.0]$), close to the instruction's 9.7. Typically, prices and units became plain text (``\textbackslash\$338'', ``87 km/h''), while one displayed equation per answer kept its \texttt{\$\$}.

\paragraph{Side effects.} On the 20 held-out how-to prompts, the edit left length unchanged (351 vs.\ 341 words; all hit the 500-token limit) and degenerate repeats unchanged (1 vs.\ 1 answer).

\subsection{E7: why self-steering broke arithmetic, read from neuron values}
\label{sec:e7}
In E4, Qwen answered arithmetic and time questions worse when it switched \texttt{lazy} on itself than when the same tool call was a placebo (Section~\ref{sec:e4}). Yet the same edit, applied without the tool context, was harmless. E7 asks why, from reads rather than outputs.

\paragraph{The failure is rare on Kaggle's weights.}
On the 12 practice questions (seed 17), Qwen in FP32 called \texttt{lazy} on 4. Those 4 answers got shorter, and 3 stayed correct; the fourth gave ``85 minutes'' where hours and minutes were asked for. That is too little failure to tune a fix on, so the diagnosis relies on reads.

\paragraph{The edit's effect flips sign with context.}
Claude read the 9 well-formed placebo answers, over the first 6 answer tokens (the \texttt{lazy} window), in four contexts. The measure is how likely Qwen finds those good answers:

\begin{center}\small
\begin{tabular}{lcc}\toprule
Mean log-prob per token & bare chat & after the tool call\\\midrule
no edit & $-1.74$ & $-1.15$\\
\texttt{lazy} on & $-1.45$ & $-1.30$\\
effect of \texttt{lazy} & $+0.30$ & $-0.15$\\\bottomrule
\end{tabular}
\end{center}

In a bare chat, \texttt{lazy} makes these openings \emph{more} likely. After the tool call it makes them less likely, on 7 of 9 answers.

This is not the overshoot we had expected, where the edit stacks on a push the tool context already makes. After the tool call, the edit's residual effect is \emph{smaller} (norm 6.7 vs.\ 9.8 at layer 29) and less aligned with the tool context's own shift (cosine 0.22 vs.\ 0.46).

\paragraph{The extra component writes digits.}
The kernel separated out the part of the edit's effect that appears only after the tool call (tool-context effect minus bare effect). Its logit lens at the last layer promotes the digits `` 9'', `` 2'', `` 3'', `` 4'' and `` 8''. It suppresses the curly quotes and punctuation that the edit promotes in a bare chat.

That part grows from layer 14 on. It is the E4 failure, read directly from the residual stream: after Qwen has announced \texttt{lazy}, the edit pushes it to start with a number.

\paragraph{Status.}
6 of 10 passes are used. Selecting which of \texttt{lazy}'s 344 neurons, or which downstream neurons, carry the digit component, and the pre-registered evaluation of a fixed edit, were cut short when the Kaggle session was cancelled. We report the diagnosis as it stands.

\subsection{E8: \texttt{eli5} from neuron values}
\label{sec:e8}
E5's \texttt{eli5} was a dense CAA vector. Here Claude had to build it as a scaling edit on chosen neurons, from reads alone, with 10 passes. It used 8 (Table~\ref{tab:e8}).

\begin{table}[t]\centering\small
\begin{tabular}{rp{6.2cm}p{5.6cm}}\toprule
Pass & What Claude ran & Result\\\midrule
1--2 & 16 practice answers, normal and with the eli5 instruction & FK 10.0 and 2.2\\
3 & read: kid-style vs.\ normal answer texts & the gap grows from L16 to L31; even exact matching of all L16--31 neurons covers 27\% of it\\
4 & set 1024 neurons to their kid-text level & no effect (FK 10.4 vs.\ 9.3 on 8 prompts)\\
5 & read: the same texts with vs.\ without the instruction & MLPs write 19\% of the instruction's effect directly, and work against it in L8--16\\
6 & 2048 neurons that already write toward it, $\times4$ & FK 6.9 vs.\ 9.3; 1/8 collapses (``money money \dots'')\\
7 & 4096 neurons ranked by alignment, L16--31, $\times4$ & FK 5.9 vs.\ 10.0; 2/16 collapse (``little little \dots'')\\
8 & the same rule without L30--31: 3584 neurons & FK 6.5 vs.\ 10.0; 0/16 collapse\\\bottomrule
\end{tabular}
\caption{E8's practice passes. FK is the Flesch--Kincaid grade.}
\label{tab:e8}
\end{table}

\paragraph{The first read was misleading.}
Contrasting a kid-style answer with a normal answer mostly measures the two texts, not the model's state. The kernel showed this before any generation was spent. If every neuron in layers 16--31 were set to its kid-text level, their direct writes would close only 27\% of the final-layer gap; the edit actually built from that read closed 12\% and did nothing.

\paragraph{A matched context fixed it.}
Pass 5 read the same texts with and without the instruction, so only the model's state differs. The logit lens of that difference at the last layer is unmistakable:
\begin{itemize}
\item it promotes `` yummy'', `` big'', `` mommy'', `` silly'' and `` sleepy'';
\item it suppresses `` fundamentally'', `` essentially'', `` predominantly'' and `` primarily''.
\end{itemize}
But MLP neurons write only 19\% of it directly, and in layers 8--16 they write against it. Attention to the instruction carries the rest.

\paragraph{So Claude stopped matching levels and started amplifying.}
A scaling edit cannot turn on a neuron that is silent without the instruction. It can only turn up neurons that already write toward the target in plain text. Claude ranked those neurons by how well their write aligns with the instruction effect.

The kernel then predicted the trade-off: every candidate edit also writes off-target, about twice as much as on-target. That off-target write showed up as collapses into repetition. Dropping the last two layers, whose neurons push words like ``little'' directly, removed the collapses on practice prompts.

\paragraph{Held-out result: a real effect, but H8 fails on degeneracy.}
On E5's 20 held-out \texttt{eli5} prompts (seeds 2000--2019, 500 tokens), \texttt{eli5-cos29} lowered the Flesch--Kincaid grade from 10.7 to 7.2. The paired difference is $-3.5$ (95\% bootstrap CI $[-4.2,-2.9]$), and the grade was lower on all 20 prompts. Answers also got shorter (274 vs.\ 336 words; 4 vs.\ 11 hit the token limit).

H8 required a drop of at least 3 points \emph{and} no degenerate outputs. One of 20 edited answers looped (``Wait, no, it is Solid State Drive (SSD).'' ten times), so H8 fails as written. Other edited answers stayed coherent but less careful. For example, ``Most vaccines do not kill you'' opens one explanation.

The instruction still did much better: grade 3.3 and 73 words, with no loops. The edit made Qwen's normal structure, headings and numbered steps, friendlier (``Here is how the magic happens'', ``like a mail delivery system''). It did not make Qwen write the short sentences of a five-year-old explanation.


